\section{Software and reproducibility}
\label{sec:software}

Every number, table and figure in this note is produced by code in the accompanying repository, and
the intention is that a reader can regenerate any of them rather than take them on trust. Nothing
in the chain reads an experiment-internal input, so ``regenerate'' means from a clone and nothing
else.

\subsection{The two modules everything rests on}

The whole budget derives from two data-free modules, and keeping them small and singular is what
makes the claim checkable:

\begin{itemize}[leftmargin=1.4em,itemsep=2pt]
  \item \code{scripts/bump\_observables.py} --- the bump observables with their fractional
resolutions, analyzable floors and published scan windows, each window carrying a per-channel
source note (Appendix~\ref{app:windows}). It also owns the two structural maps that decide what
counts as one spectrum: the same-axis merges and the lepton-flavour split. It is imported, never
duplicated; a floor or a window redefined in a consumer would be a bug.
  \item \code{scripts/public\_obs\_map.py} --- the map from public BSM model classes onto the spectra
they populate, plus the published event-selection multiplicity per spectrum.
\end{itemize}

Both carry import-time assertions that keep the flavour layer complete: every leaf channel must have
a window and a floor, and no flavour-inclusive parent may keep one, which is what prevents a channel
being counted both as itself and as part of its parent.

\subsection{Layout and rebuilding}

\begin{itemize}[leftmargin=1.4em,itemsep=2pt]
  \item \code{scripts/}, the pipeline, one concern per file.
  \item \code{results/tables/}, machine-readable output, including the cached Monte Carlo of
Section~\ref{sec:selection}.
  \item \code{results/plots/}, the figures of this note.
  \item \code{results/overviews/}, the long-form reports the sections here condense.
  \item \code{docs/}, holding \code{METHOD\_NOTES.md} (conventions and the reasoning behind the
counting choices), \code{OUTPUTS.md} (every output file mapped to the script that writes it), and
this note.
\end{itemize}

The pipeline is driven by a single makefile whose targets are real files with real prerequisites,
so a rerun redoes only what is stale:

\begin{verbatim}
    pip install -r requirements.txt   # numpy, scipy, matplotlib
    make all                          # budget + ab + stats
    make help                         # the individual stages
\end{verbatim}

The stages are \code{budget} (Sections~\ref{sec:budget} to~\ref{sec:excess}), \code{ab}
(Section~\ref{sec:ab}) and \code{stats} (Section~\ref{sec:selection}). The budget itself uses only
the Python standard library; \code{ab} and \code{stats} need NumPy~\cite{numpy},
SciPy~\cite{scipy} and Matplotlib~\cite{matplotlib}.

\subsection{Determinism}

All Monte Carlo is explicitly seeded, and rebuilding the figure set reproduces the committed files
byte for byte. That is worth more than it sounds, because it means a changed figure signals that an
input or a method changed rather than that a random seed moved. The two genuinely slow computations,
the Benjamini--Hochberg $q$ scan and its imperfect-estimator variant, are cached as \code{.npz}
files and rerun only on request.

\subsection{Relation to the catalogue note}

One question is deliberately absent here: which of these 46 spectra a given collaboration has
actually produced signal samples for. That is a statement about a production program rather than
about the search budget, it cannot be made from public information, and answering it does not change
any number in this note --- the windows are published-search windows either way. It is treated in a
companion ATLAS-internal note, which imports the two modules above unchanged so that the two
documents cannot drift apart on what a spectrum is.
